\documentclass[10pt]{article}

\usepackage[utf8]{inputenc}
\usepackage[T1]{fontenc}
\usepackage{amsmath,amssymb}
\usepackage{booktabs}
\usepackage{enumitem}
\usepackage{geometry}
\usepackage{hyperref}

\geometry{margin=0.8in}
\hypersetup{colorlinks=true,linkcolor=blue,urlcolor=blue}
\setlist{itemsep=2pt,topsep=3pt,parsep=0pt,partopsep=0pt}
\setlength{\parindent}{0pt}
\setlength{\parskip}{3pt}
\setlength{\abovedisplayskip}{5pt plus 2pt minus 2pt}
\setlength{\belowdisplayskip}{5pt plus 2pt minus 2pt}

\newcommand{\R}{\mathbb{R}}
\newcommand{\E}{\mathbb{E}}
\newcommand{\Var}{\operatorname{Var}}
\newcommand{\Cov}{\operatorname{Cov}}
\newcommand{\rank}{\operatorname{rank}}
\newcommand{\im}{\operatorname{Im}}
\newcommand{\tr}{\operatorname{tr}}
\newcommand{\diag}{\operatorname{diag}}
\newcommand{\Span}{\operatorname{span}}

\title{Summer Bootcamp -- Session 1 Exercises\\Linear Algebra}
\author{Nathan Sauldubois}
\date{August 2026}

\begin{document}
\maketitle

\section*{Instructions}

The goal is to become operational with matrix language. Exercises 1--6 are the
core paper exercises. Exercises 7--8 are hands-on Python exercises. Exercises
9--11 are decomposition brainteasers and applications. Exercises 12--13 are
finance-flavored brainteasers and theoretical questions. Exercises 14--16
connect transition matrices with Markov chains and hitting-time puzzles.

For the hands-on exercises, use only basic \texttt{NumPy}; do not use a built-in
PCA or portfolio-optimization function unless explicitly asked to compare.

\section{Core Linear Algebra}

\paragraph{Exercise 1 -- Matrix arithmetic and dimension checks.}
Let
\[
A=\begin{pmatrix}
1&2&-1\\
0&3&4
\end{pmatrix},\qquad
B=\begin{pmatrix}
2&1\\
-1&0\\
3&2
\end{pmatrix},\qquad
C=\begin{pmatrix}
2&-1\\
1&1
\end{pmatrix},
\]
and
\[
x=\begin{pmatrix}1\\-2\\0\end{pmatrix},\qquad
y=\begin{pmatrix}3\\1\end{pmatrix}.
\]

\begin{enumerate}
\item For each expression, say whether it is defined and give its dimension:
\[
AB,\quad BA,\quad AC,\quad CA,\quad Ax,\quad Ay,\quad Bx,\quad By,\quad
C^2,\quad A^\top C.
\]
\item Compute all defined products among $AB$, $BA$, $Ax$, $By$, and $C^2$.
\item Compute $A^\top$, $B^\top$, $\tr(C)$, $\det(C)$, and $C^{-1}$.
\item Use your computations to show explicitly that matrix multiplication is
not commutative.
\item Interpret $Ax$ as a linear combination of the columns of $A$.
\end{enumerate}

\paragraph{Exercise 2 -- Kernel, image, rank, and geometry.}
Consider
\[
L=
\begin{pmatrix}
1&-1&0\\
-1&2&-1\\
0&-1&1
\end{pmatrix}.
\]

\begin{enumerate}
\item Solve $Lx=0$ and give a basis of $\ker(L)$.
\item Compute $\rank(L)$ using rank-nullity.
\item Show that every column of $L$ is orthogonal to $(1,1,1)^\top$.
\item Deduce that
\[
\im(L)=\{y\in\R^3:y_1+y_2+y_3=0\}.
\]
\item For each vector below, decide whether $Lx=b$ has a solution:
\[
b_1=\begin{pmatrix}1\\0\\-1\end{pmatrix},\qquad
b_2=\begin{pmatrix}1\\1\\0\end{pmatrix}.
\]
When a solution exists, find all solutions.
\item Give a geometric interpretation: which direction is killed by $L$, and
on which plane does $L$ act non-trivially?
\end{enumerate}

\paragraph{Exercise 3 -- Linear systems and bootstrapping.}

\textbf{Part A: a small system.}
Consider
\[
\begin{cases}
2x_1+3x_2=5,\\
x_1-4x_2=-2.
\end{cases}
\]
\begin{enumerate}
\item Write the system as $Ax=b$.
\item Compute $\det(A)$ and decide whether the solution is unique.
\item Solve the system by substitution.
\item Compute $A^{-1}$ and verify that $x=A^{-1}b$ gives the same answer.
\end{enumerate}

\textbf{Part B: fixed-income bootstrapping.}
You observe three annual coupon bonds with face value $1$:
\[
\begin{array}{@{}cccc@{}}
\toprule
\text{Maturity} & \text{Coupon} & \text{Cash flows} & \text{Price}\\
\midrule
1 & 3\% & 1.03\text{ at }T_1 & 1.005\\
2 & 4\% & 0.04\text{ at }T_1,\ 1.04\text{ at }T_2 & 1.018\\
3 & 5\% & 0.05\text{ at }T_1,\ 0.05\text{ at }T_2,\ 1.05\text{ at }T_3 & 1.025\\
\bottomrule
\end{array}
\]
Let $d_k=D(T_k)$ be the discount factor at maturity $T_k$.
\begin{enumerate}
\item Write the pricing equations in matrix form $Cd=p$.
\item Explain why the system is triangular.
\item Solve for $d_1,d_2,d_3$ using forward substitution.
\item Compute the zero rates
\[
z_k=-\frac{1}{T_k}\log d_k,\qquad k=1,2,3.
\]
\item Explain why this is a linear algebra version of bootstrapping.
\end{enumerate}

\paragraph{Exercise 4 -- Covariance matrices and portfolio variance.}
The following centered return matrix contains four daily observations of two
assets:
\[
X=
\begin{pmatrix}
0.01&-0.02\\
0.02&0.01\\
-0.01&0.00\\
-0.02&0.01
\end{pmatrix}.
\]

\begin{enumerate}
\item Check that each column is centered.
\item Compute
\[
\widehat{\Sigma}=\frac{1}{n-1}X^\top X,\qquad n=4.
\]
\item Verify directly that $\widehat{\Sigma}$ is symmetric.
\item For $w=(0.6,0.4)^\top$, compute $w^\top\widehat{\Sigma}w$.
\item Decompose the variance into individual variance terms and covariance
terms:
\[
w^\top\widehat{\Sigma}w
=w_1^2\sigma_1^2+w_2^2\sigma_2^2+2w_1w_2\sigma_{12}.
\]
\item Explain why covariance can either increase or reduce portfolio risk.
\end{enumerate}

\paragraph{Exercise 5 -- Diagonalization and covariance interpretation.}
For $\rho\in(-1,1)$, consider the two-asset correlation matrix
\[
C_\rho=
\begin{pmatrix}
1&\rho\\
\rho&1
\end{pmatrix}.
\]

\begin{enumerate}
\item Compute the eigenvalues and eigenvectors of $C_\rho$.
\item Specialize to $\rho=0.8$ and write an orthogonal diagonalization
$C_\rho=QDQ^\top$.
\item Interpret the two eigenvectors as the common direction and the spread
direction.
\item Which eigenvalue becomes small when $\rho$ is close to $1$? What does
that mean financially?
\item What happens when $\rho=-1$? Is the matrix invertible?
\end{enumerate}

\paragraph{Exercise 6 -- PCA by hand.}
The centered data matrix is
\[
X=
\begin{pmatrix}
-1&-1\\
0&0\\
1&1
\end{pmatrix}.
\]

\begin{enumerate}
\item Compute $\widehat{\Sigma}=\frac{1}{n-1}X^\top X$.
\item Find its eigenvalues and orthonormal eigenvectors.
\item Identify the first principal component.
\item Project the three observations onto the first principal component.
\item Reconstruct the observations using only the first principal component.
Is the reconstruction exact?
\item Explain what information is kept and what information is lost.
\end{enumerate}

\section{Hands-On Python}

\paragraph{Exercise 7 -- Empirical covariance and PCA with NumPy.}
Use the following return matrix. Rows are days and columns are assets:
\[
R=
\begin{pmatrix}
0.010&-0.004&0.006&0.012\\
-0.006&0.003&-0.002&-0.008\\
0.014&0.001&0.010&0.015\\
-0.012&-0.006&-0.007&-0.010\\
0.008&0.004&0.005&0.009\\
0.003&-0.002&0.001&0.004\\
-0.009&0.002&-0.004&-0.011\\
0.011&0.005&0.007&0.013
\end{pmatrix}.
\]

\begin{enumerate}
\item Load this matrix into a \texttt{NumPy} array.
\item Compute the empirical mean vector.
\item Center the matrix manually.
\item Compute the empirical covariance matrix using
\[
\widehat{\Sigma}=\frac{1}{n-1}R_c^\top R_c.
\]
\item Compare your result with \texttt{np.cov(R, rowvar=False)}.
\item Compute eigenvalues and eigenvectors using \texttt{np.linalg.eigh}.
\item Sort eigenvalues in decreasing order and compute the explained variance
ratios.
\item Project the centered data onto the first one and then the first two
principal components.
\item In two or three sentences, interpret the first principal component as a
financial factor.
\end{enumerate}

Suggested skeleton:
\begin{verbatim}
import numpy as np

R = np.array([
    [ 0.010, -0.004,  0.006,  0.012],
    [-0.006,  0.003, -0.002, -0.008],
    [ 0.014,  0.001,  0.010,  0.015],
    [-0.012, -0.006, -0.007, -0.010],
    [ 0.008,  0.004,  0.005,  0.009],
    [ 0.003, -0.002,  0.001,  0.004],
    [-0.009,  0.002, -0.004, -0.011],
    [ 0.011,  0.005,  0.007,  0.013],
])

n, d = R.shape
mu = ...
Rc = ...
Sigma = ...
eigvals, eigvecs = ...
\end{verbatim}

\paragraph{Exercise 8 -- Portfolio risk lab.}
Consider four assets with expected returns and covariance matrix
\[
\mu=
\begin{pmatrix}
0.05\\0.07\\0.09\\0.12
\end{pmatrix},
\qquad
\Sigma=
\begin{pmatrix}
0.040&0.012&0.018&0.010\\
0.012&0.060&0.020&0.016\\
0.018&0.020&0.090&0.030\\
0.010&0.016&0.030&0.160
\end{pmatrix}.
\]

\begin{enumerate}
\item Check numerically that $\Sigma$ is symmetric.
\item Check numerically that all eigenvalues of $\Sigma$ are positive.
\item For
\[
w_0=(0.25,0.25,0.25,0.25)^\top,
\]
compute expected return, variance, and volatility.
\item Compute the unconstrained minimum-variance portfolio with budget
$\mathbf{1}^\top w=1$:
\[
w^\star=\frac{\Sigma^{-1}\mathbf{1}}
{\mathbf{1}^\top\Sigma^{-1}\mathbf{1}}.
\]
Use \texttt{np.linalg.solve} rather than explicitly forming the inverse.
\item Compare $w^\star$ to $w_0$ in terms of return and volatility.
\item Generate 5000 random long-only portfolios, normalize each one so weights
sum to $1$, and plot or print the lowest volatility you find.
\item Does the best random portfolio get close to $w^\star$? Explain why or
why not.
\end{enumerate}

\section{Decomposition Brainteasers and Applications}

\paragraph{Exercise 9 -- Predicting a coupled system without iterating.}
Two quantities evolve according to
\[
x_{k+1}=Ax_k,
\qquad
A=\begin{pmatrix}3&1\\0&2\end{pmatrix},
\qquad
x_0=\begin{pmatrix}0\\1\end{pmatrix}.
\]

\begin{enumerate}
\item Compute $x_1,x_2,x_3$ directly and guess the form of $x_k$.
\item Diagonalize $A$ and use the decomposition to compute $A^k$.
\item Deduce a closed formula for $x_k$ and verify your initial guess.
\item Which eigenvalue controls the long-run growth? Explain without carrying
out any additional matrix multiplication.
\item Suppose instead that $x_0=(-1,1)^\top$. What simplification occurs, and
why?
\end{enumerate}

\paragraph{Exercise 10 -- Finding hidden risk factors.}
Two standardized asset returns have correlation matrix
\[
C=\begin{pmatrix}1&0.8\\0.8&1\end{pmatrix}.
\]

\begin{enumerate}
\item Without expanding a determinant, propose two orthogonal candidate
eigenvectors using the symmetry between the assets.
\item Normalize these vectors and compute their eigenvalues.
\item Interpret the two eigenportfolios as a common trade and a spread trade.
\item Compute the variance of each unit eigenportfolio.
\item A trader claims that the equally weighted long--short portfolio is almost
riskless. Is that statement literally correct? Quantify your answer.
\item Repeat the reasoning for
$C_\rho=\begin{pmatrix}1&\rho\\\rho&1\end{pmatrix}$ and describe what happens
as $\rho\to1$ and as $\rho\to-1$.
\end{enumerate}

\paragraph{Exercise 11 -- Manufacturing a target correlation.}
Let $Z_1,Z_2$ be independent standard Gaussian variables. You want to construct
a Gaussian vector $X=(X_1,X_2)^\top$ with covariance matrix
\[
\Sigma=\begin{pmatrix}4&2\\2&3\end{pmatrix}.
\]

\begin{enumerate}
\item Find the lower-triangular Cholesky factor $L$ such that
$\Sigma=LL^\top$.
\item Write $X=LZ$ component by component.
\item Compute $\Var(X_1)$, $\Var(X_2)$, and $\Cov(X_1,X_2)$ directly from
your formulas and verify that they match $\Sigma$.
\item Compute the correlation between $X_1$ and $X_2$.
\item If $z=(1,-1)^\top$ is one simulated independent shock, compute the
corresponding correlated shock $Lz$.
\item Explain why multiplying by $\Sigma$ itself, instead of $L$, would not
produce covariance $\Sigma$.
\end{enumerate}

\section{Finance Brainteasers and Theory}

\paragraph{Exercise 12 -- Singular covariance and redundant assets.}
Let $R_1$ and $R_2$ be two returns with
\[
\Var(R_1)=1,\qquad \Var(R_2)=1,\qquad \Cov(R_1,R_2)=0.2.
\]
Define a third return by $R_3=R_1+R_2$.

\begin{enumerate}
\item Compute $\Cov(R_1,R_3)$, $\Cov(R_2,R_3)$, and $\Var(R_3)$ from the
identity $R_3=R_1+R_2$.
\item Use these quantities to write the full covariance matrix $\Sigma$ of
$(R_1,R_2,R_3)$.
\item Write $R=(R_1,R_2,R_3)^\top$ as $R=MX$, where
$X=(R_1,R_2)^\top$, and deduce the factorization
$\Sigma=M\Sigma_XM^\top$.
\item Find a nonzero vector $w$ such that $\Sigma w=0$ and verify the matrix
multiplication explicitly.
\item Compute and interpret the corresponding portfolio return $w^\top R$.
\item Explain in three equivalent ways why $\Sigma$ is singular: using the
linear relation among returns, the kernel, and the rank of $M$.
\item Suppose three securities have payoffs exactly $R_1,R_2,R_3$ and current
prices $p_1,p_2,p_3$. Derive the price relation required by no arbitrage.
\item If $p_3<p_1+p_2$, give the positions, initial cash flow, and terminal
payoff. Repeat the analysis for $p_3>p_1+p_2$.
\end{enumerate}

\paragraph{Exercise 13 -- Short theoretical questions.}
Give short but rigorous answers.

\begin{enumerate}
\item Prove that for any real matrix $X$, the matrix $X^\top X$ is symmetric
positive semidefinite.
\item Prove that if $A$ is injective, then $\ker(A)=\{0\}$, and prove the
converse.
\item Let $A\in\R^{m\times n}$. Explain why $Ax=b$ has a solution if and only
if $b\in\im(A)$.
\item Let $A$ be real symmetric. Prove that eigenvectors associated with
distinct eigenvalues are orthogonal.
\item Give an example of a matrix with one repeated eigenvalue that is not
diagonalizable.
\item Explain, in one paragraph, why PCA is just diagonalization of a covariance
matrix.
\end{enumerate}

\section{Markov Chains and Matrix Brainteasers}

Throughout this section, probability distributions are column vectors and
$P_{ij}=\mathbb{P}(X_{n+1}=i\mid X_n=j)$.

\paragraph{Exercise 14 -- Credit migration by matrix powers.}
Consider
\[
P=\begin{pmatrix}0.9&0.4\\0.1&0.6\end{pmatrix}.
\]

\begin{enumerate}
\item Verify that $P$ is column-stochastic and explain the meaning of each
column.
\item Starting from $\mu_0=(1,0)^\top$, compute $\mu_1$ and $\mu_2$.
\item Find the stationary distribution $\pi$ by solving
$P\pi=\pi$ together with $\mathbf{1}^\top\pi=1$.
\item Find both eigenvalues of $P$ and diagonalize it.
\item Obtain a closed formula for $\mu_n=P^n\mu_0$.
\item How many periods are needed before the Bad-state probability is within
$10^{-3}$ of its stationary value?
\end{enumerate}

\paragraph{Exercise 15 -- Waiting time for PPF.}
A fair coin is tossed repeatedly until the consecutive pattern PPF appears,
where P denotes heads and F denotes tails.

\begin{enumerate}
\item Define states according to the longest suffix that is also a prefix of
PPF.
\item Draw the transition graph and write its transition matrix, including an
absorbing completion state.
\item Let $E_0,E_P,E_{PP}$ be the expected remaining waiting times. Condition
on the next toss to obtain a linear system.
\item Solve the system and compute the expected number of tosses from the
start.
\item Explain why observing another P from state PP does not reset the process.
\end{enumerate}

\paragraph{Exercise 16 -- Waiting time for 564.}
A fair six-sided die is rolled repeatedly until three consecutive rolls are
5, then 6, then 4.

\begin{enumerate}
\item Construct states $0$, $5$, and $56$ using the longest useful suffix.
\item Write the transition matrix after adding an absorbing state for 564.
\item Let $E_0,E_5,E_{56}$ be the expected remaining waiting times. Derive the
three first-step recurrences.
\item Solve the resulting linear system.
\item Compare the answer with $6^3$ and explain the relationship.
\end{enumerate}

\end{document}
